\documentclass[../../../include/open-logic-section]{subfiles}

\begin{document}

\olfileid[pt]{his}{set}{infinitesimals}
\olsection{Infinitesimais e Diferenciação}

Newton e Leibniz descobriram o cálculo (independentemente) no fim do
século XVII. Uma aplicação particularmente importante do cálculo era a
\emph{diferenciação}. Grosso modo, a diferenciação
pretende dar uma noção da ``taxa de variação'', ou gradiente, de uma
função num ponto.

Eis uma forma vívida de ilustrar a ideia. Considere a função
$f(x) = \nicefrac{x^2}{4} + \nicefrac{1}{2}$, representada a preto abaixo:
\begin{center}
	\begin{tikzpicture}[scale=1]
	\draw[->, gray] (-1,0) -- (4.25,0) node[right] {$x$};
	\draw[->, gray] (0,-0.25) -- (0,5.25) node[above] {$f(x)$};
	
	\foreach \x/\xtext in {1/1, 2/2, 3/3, 4/4}
	\draw[shift={(\x,0)}, gray] (0pt,2pt) -- (0pt,-2pt) node[below] {$\xtext$};
	
	\foreach \y/\ytext in {1/1, 2/2, 3/3, 4/4, 5/5}
	\draw[shift={(0,\y)}, gray] (2pt,0pt) -- (-2pt,0pt) node[left] {$\ytext$};
	
	\draw[oldiagcolorC] (0.5, 9/16) -- (3.5, 57/16) -- (3.5, 9/16)--cycle;
	\draw[oldiagcolorD] (.5, 9/16) -- (2.5, 33/16) -- (2.5, 9/16) -- cycle;
	\draw[oldiagcolorE] (.5, 9/16) -- (1.5, 17/16) -- (1.5, 9/16) -- cycle;
	\draw[thick] (-1,.75) parabola bend (0,0.5) (4,4.5);
	\end{tikzpicture}
\end{center}
Suponhamos que queremos encontrar o gradiente da função em $c =
\nicefrac{1}{2}$. Começamos por desenhar um triângulo cuja hipotenusa
aproxima o gradiente nesse ponto, talvez o triângulo vermelho acima.
Quando $\beta$ é o comprimento da base do nosso triângulo, a sua altura é
$f(\nicefrac{1}{2}+\beta) - f(\nicefrac{1}{2})$, de modo que o gradiente
da hipotenusa é:
\[
\frac{f(\nicefrac{1}{2}+\beta) - f(\nicefrac{1}{2})}{\beta}.
\]
Assim, o gradiente do nosso triângulo !!{colorC}, com base de comprimento~$3$, é
exatamente~$1$. A hipotenusa de um triângulo menor, o triângulo !!{colorD}
com base de comprimento~$2$, dá uma melhor aproximação; o seu
gradiente é $\nicefrac{3}{4}$. Um triângulo ainda menor, o !!{colorE}
com base de comprimento~$1$, dá uma aproximação ainda melhor, com
gradiente $\nicefrac{1}{2}$. 

Triângulos cada vez menores dão-nos aproximações cada vez melhores. Por isso
poderíamos dizer algo deste gênero: a hipotenusa de um triângulo com
comprimento de base \emph{infinitesimal} nos dá o próprio gradiente em $c =
\nicefrac{1}{2}$. Deste modo, obteríamos uma fórmula para a
(primeira) derivada da função $f$ no ponto $c$:
\[
{f'}(c) = \frac{f(c+\beta) - f(c)}{\beta} \text{ onde $\beta$ é infinitesimal.}
\]
E, grosso modo, foi isto que Newton e Leibniz disseram.

Contudo, uma vez que o disseram, devemos perguntar-lhes: o que é um
\emph{infinitesimal}? Surge um dilema sério. Se $\beta = 0$, então
$f'$ fica mal definida, pois envolve dividir por $0$. Mas se $\beta >
0$, obtemos apenas uma \emph{aproximação} ao gradiente, e não o
próprio gradiente.

Esta não é uma preocupação anacrônica. Eis Berkeley, a criticar os
seguidores de Newton:
\begin{quote}
Admito que os signos possam ser feitos para denotar qualquer coisa ou nada:
e consequentemente que, na notação original $c + \beta$, $\beta$
poderia ter significado um incremento ou nada. Mas então, seja qual for
o significado atribuído a $\beta$, é preciso argumentar de modo consistente com
essa significação, e não proceder com duplo sentido: o que seria um sofismo manifesto. (\citealt[\S{}XIII]{Berkeley1734},
variáveis alteradas para coincidir com o texto precedente)
\end{quote}
Para defender o cálculo infinitesimal face a Berkeley, poder-se-ia responder
que a linguagem dos ``infinitesimais'' é meramente figurada. Poder-se-ia
dizer que, desde que tomemos um triângulo \emph{bem pequeno}, obteremos
uma aproximação \emph{suficientemente boa} à tangente. Berkeley tinha
também uma resposta para isto: embora isso pudesse bastar para a engenharia,
minava o \emph{estatuto} da matemática, pois
\begin{quote}
dizem-nos que \emph{in rebus mathematicis errores qu\`{a}m minimi
non sunt contemnendi}. [No caso da matemática, os menores
erros não devem ser desprezados.] \citep[\S{}IX]{Berkeley1734}
\end{quote}
A passagem em itálico é uma citação quase literal do próprio Newton na
\emph{Quadrature of Curves} (1704).

As objeções filosóficas de Berkeley são profundamente incisivas. Não obstante,
o cálculo foi uma empresa enormemente bem-sucedida, e os matemáticos
continuaram a usá-lo sem incorrer em erro.

\end{document}
